% ===================================================================
%  CHART No. 2 — The Human Body. — Break. — Games.
%  Traduction 2026 d'apres texto/io, controlee sur texto/fr.
%  Consigne : voir texto/en/10-table-01.tex.
% ===================================================================

%%K t02-apar-1 apar
\VUsaut{4.0mm}
\VUtitre{15.0pt}{60}{CHART No. 2}
\VUsaut{2.0mm}
\VUtitre{11.4pt}{0}{The Human Body. --- Break.}
\VUtitre{11.4pt}{0}{Games.}

%%K t02-tit-0 sub
\VUtitre{13.2pt}{0}{The Human Body.}
\VUsaut{2.4mm}
\VUcentre{10.2pt}{0}{\textit{(A simple natural science lesson.)}}
\VUsaut{3.0mm}

%%K t02-01-1 p
1. --- The main parts of the human body are: the \VUgras{head}, the
\VUgras{trunk} and the \VUgras{limbs}.

%%K t02-tit-1 sub
\VUpk{10.2pt}{40}{The Head.}
\VUsaut{3.0mm}

%%K t02-02-1 p
2. --- The head has two parts: the \VUgras{skull}, which holds the
\VUgras{brain} and is covered by the \VUgras{hair}, and the \VUgras{face}.

%%K t02-03-1 p
3. --- In our drawing we can see neither the skull nor the brain; but we can see the important
parts of the face very clearly.

%%K t02-04-1 p
4. --- First comes the \VUgras{forehead}, which suggests a powerful intellect and a mind
always at work. The \VUgras{temples} are very open. The \VUgras{eye} is lively and wide open.
A thick \VUgras{eyebrow} and long
\VUgras{eyelashes} shield it.

%%K t02-05-1 p
5. --- Then come the \VUgras{nose} and the \VUgras{nostrils}. We use the nose to smell and to
breathe. On either side of the nose are the \VUgras{cheeks}, and further out the
\VUgras{ears}. The lower part of the face consists of the \VUgras{mouth} and the
\VUgras{chin}.

%%K t02-06-1 p
6. --- We cannot see them very well, because the \VUgras{moustache} and the
\VUgras{beard} hide them from us. But we know that the mouth has the two
\VUgras{lips}, the \VUgras{upper jaw} and the \VUgras{lower jaw} with their
\VUgras{teeth}, and the \VUgras{tongue}. With the mouth we speak and eat.
With the tongue we taste our food.

%%K t02-07-1 p
7. --- The eye, the ear, the nose and the mouth are, together with the fingers, the organs of
the five senses: \VUgras{sight}, \VUgras{hearing},
\VUgras{smell}, \VUgras{taste} and \VUgras{touch}.

%%K t02-08-1 p
8. --- The head is therefore one of the most interesting parts of the human body.

%%K t02-tit-2 sub
\VUsaut{4.4mm}
\VUpk{10.2pt}{40}{The Trunk.}
\VUsaut{3.0mm}

%%K t02-09-1 p
9. --- The \VUgras{neck} joins the head to the trunk. It includes the
\VUgras{throat} and the \VUgras{nape}.

%%K t02-10-1 p
10. --- The trunk also holds many essential organs. In the \VUgras{chest} are the
\VUgras{lungs}, with which we breathe, and the \VUgras{heart}, which is the centre of life.
When the heart stops beating, the living creature dies.

%%K t02-11-1 p
11. --- The \VUgras{ribs} support the chest.

%%K t02-12-1 p
12. --- The other parts of the trunk are the \VUgras{side}, the
\VUgras{back}, the \VUgras{shoulders} and the \VUgras{abdomen} (belly). We lie
on our side or on our back to sleep, and we carry loads on our shoulder.

%%K t02-13-1 p
13. --- In the abdomen are the \VUgras{stomach} and the \VUgras{intestines}. They are what
digests our food.

%%K t02-tit-3 sub
\VUsaut{4.4mm}
\VUpk{10.2pt}{40}{The Limbs.}
\VUsaut{3.0mm}

%%K t02-14-1 p
14. --- We have four \VUgras{limbs}: two \VUgras{arms} and two \VUgras{legs}. With our arms we
take, hold, lift and gather up objects; with our legs we stand, walk and run.

%%K t02-15-1 p
15. --- The \VUgras{shoulder} joins the arm to the body. A very sturdy
\VUgras{muscle}, the biceps, moves the arm, which the \VUgras{elbow} joins to
the \VUgras{forearm}. The \VUgras{wrist} forms the joint of the \VUgras{hand}, which ends in
the \VUgras{fingers} and the \VUgras{nails}. On the hand we can make out a \VUgras{vein} (and
an artery) in which the \VUgras{blood} flows.

%%K t02-16-1 p
16. --- The parts of the leg are: the \VUgras{thigh}, the \VUgras{knee} and the
\VUgras{back of the knee}, the \VUgras{calf}, whose \VUgras{flesh} is covered by the
\VUgras{skin}, the \VUgras{ankle} and the \VUgras{foot}. The
\VUgras{bones} of the leg are very thick and very strong. At the end of the
foot are the \VUgras{toes}. The other parts are the \VUgras{sole} and the
\VUgras{heel}.

%%K t02-17-1 p
17. --- That is a nearly complete description of the human body.

%%K t02-17-2 p
\textit{Note.} --- \textit{With the help of the phrase “my ... hurts (my head, and so on)”,
the teacher can go over this whole vocabulary again from the point of view of good or bad}
\VUgras{bodily health}.

%%K t02-tit-4 sub
\VUsaut{5.0mm}
\VUpk{10.2pt}{40}{Gymnastics.}
\VUsaut{2.4mm}
\VUcentre{10.2pt}{0}{\textit{(Told by one of the pupils.)}}
\VUsaut{3.0mm}

%%K t02-18-1 p
18. --- To be supple, nimble and healthy, we must exercise our legs and our arms. That is why
we play games and do \VUgras{gymnastics}.

%%K t02-19-1 p
19. --- During the week these two exercises take place in the school yard (see chart No. 1).
But on Thursdays and Sundays we go out to a large lawn, where we have room to run and enjoy
ourselves.

%%K t02-20-1 p
20. --- Our teachers and our parents sometimes come with us; but it is the
\VUgras{gymnastics teacher} \textsuperscript{(12)} who directs the exercises.

%%K t02-21-1 p
21. --- On the lawn, to the left of the entrance, stands the
\VUgras{gymnastic apparatus} \textsuperscript{(1)}: we climb the \VUgras{ladder}
\textsuperscript{(2)}, we swing upside down on the \VUgras{rings} \textsuperscript{(3)} and on
the \VUgras{trapeze} \textsuperscript{(4)}, we haul ourselves up hand over hand to the top of
the \VUgras{rope ladder} \textsuperscript{(5)} or of the
\VUgras{knotted rope} \textsuperscript{(6)}.

%%K t02-22-1 p
22. --- Meanwhile our friends are vaulting over the \VUgras{wooden horse}
\textsuperscript{(7)} or the \VUgras{parallel bars} \textsuperscript{(11)}. Others are
\VUgras{boxing} \textsuperscript{(8)} or \VUgras{fencing} \textsuperscript{(9)} with a mask
and a \VUgras{foil}. Others, finally, are lifting \VUgras{dumbbells} \textsuperscript{(10)}.

%%K t02-tit-5 sub
\VUsaut{4.6mm}
\VUpk{10.2pt}{40}{Games.}
\VUsaut{3.0mm}

%%K t02-23-1 p
23. --- Around the gymnasts, boys and girls are playing various
\VUgras{games}. The girls are playing with their \VUgras{hoop}
\textsuperscript{(14)}; they \VUgras{skip} \textsuperscript{(13)}, or invite their brothers
and cousins to a game of \VUgras{croquet} \textsuperscript{(15)}. They love knocking their
opponent's \VUgras{ball} away with their
\VUgras{wooden mallet}, just when he thinks he is going straight through the
hoop!

%%K t02-24-1 p
24. --- The boys prefer more strenuous games. They gladly play
\VUgras{skittles} \textsuperscript{(16)}, which they knock down neatly with the
\VUgras{ball} \textsuperscript{(17)}. They do not turn up their noses at the \VUgras{spinning top}
\textsuperscript{(18)} and \VUgras{cup-and-ball} \textsuperscript{(19)}, or even
\VUgras{toad in the hole} \textsuperscript{(20)}. But their favourite games are
\VUgras{marbles} \textsuperscript{(21)} and \VUgras{fives} \textsuperscript{(22)}, because the
wall of the \VUgras{greenhouse} \textsuperscript{(26)} is just right for bouncing the light
ball off.

%%K t02-25-1 p
25. --- Little John, up on his \VUgras{stilts} \textsuperscript{(24)}, is very skilful and
keeps his balance very well. Near him, a few friends are playing \VUgras{hide-and-seek}
\textsuperscript{(25)} in that greenhouse and behind the \VUgras{hedge}. Others prefer
\VUgras{hot cockles} \textsuperscript{(28)} or
\VUgras{shuttlecock} \textsuperscript{(29)}, which flits from one \VUgras{racket} to
the other.

%%K t02-noto-1 noto
\VUnotes{143.2mm}{%
(*) Children's games often have purely national names. We have named them in Ido as best we
could, either by drawing on the very action that characterises them, or by adopting the
national name of one country or another when it is not too misleading. For example: the
\VUgras{barrel game} (German and French) is the \VUgras{toad in the hole}
\textsuperscript{(20)} (English), in which the players try to throw their round discs through
the mouth of the metal toad that stands gaping on the pierced box (barrel).
\VUgras{Shoulder-leap} \textsuperscript{(30)} differs from
\VUgras{leapfrog} only in this: to jump over the head, the players rest
their hands on their partner's shoulders, whereas in “\VUgras{leapfrog}” they rest their hands
only on his back. --- Or again, some of the players bend over while the others vault over
their backs with their hands on their shoulders; the first must take the shock without giving
way, the second must jump without losing their balance (see the chart itself).%
}

%%K t02-26-1 p
26. --- In the middle of the lawn there are two trees. At the foot of one of them, seven or
eight boys have started a game of \VUgras{shoulder-leap} \textsuperscript{(30)} (*). This game
is not known in every country; but everyone knows what \VUgras{leapfrog} is, which the
children are not playing just now.

%%K t02-tit-6 sub
\VUsaut{5.0mm}
\VUpk{10.2pt}{40}{Games in the open air.}
\VUsaut{3.4mm}

%%K t02-27-1 p
27. --- \VUgras{Blind man's buff} \textsuperscript{(31)} is also great fun; girls and boys
take turns to put the \VUgras{blindfold} over their eyes and catch the clumsy ones who let
themselves be caught.

%%K t02-28-1 p
28. --- In the middle of the lawn, here is a very French game, though a somewhat rough one: it
is the \VUgras{bear} \textsuperscript{(32)}, far noisier than the peaceful game of the
\VUgras{kite} \textsuperscript{(33)} or even than the English game of \VUgras{tennis}
\textsuperscript{(34)}.

%%K t02-29-1 p
29. --- That last sport is what the older pupils love best. Even our teachers are happy to
join in. It calls for real skill and agility, and I love it. I leave the \VUgras{swing}
\textsuperscript{(35)} to the girls.

%%K t02-30-1 p
30. --- I am not keen on another English game that could easily turn dangerous.

%%K t02-30-1b p
I mean \VUgras{football} \textsuperscript{(36)}, which my older brother loves. I prefer our
old game of \VUgras{prisoner's base} \textsuperscript{(38)}, just as healthy and far less
brutal.

%%K t02-tit-7 sub
\VUsaut{4.6mm}
\VUpk{10.2pt}{40}{Cycling and ball games.}
\VUsaut{3.0mm}

%%K t02-31-1 p
31. --- I also love riding round the lawn on a \VUgras{bicycle} \textsuperscript{(37)}.

%%K t02-32-1 p
32. --- Next year I am going to learn to \VUgras{ride} \textsuperscript{(39)}. A horse is such
a fine sight when it steps or trots gracefully, and clears the obstacles at a gallop!

%%K t02-33-1 p
33. --- In summer we learn to \VUgras{swim} \textsuperscript{(40)}. We dive and bathe with
delight in the clear, cool water of the river. Sometimes, at the end, we send up a paper
\VUgras{balloon} \textsuperscript{(41)}, which rises majestically into the air as soon as it
is filled with hot air.

%%K t02-34-1 p
34. --- There are many other games as well, but I would never get to the end of the list, and
you can see from this short account that Thursdays on our lawn are anything but dull.

%%K t02-34-2 p
N.B. --- \textit{The teacher will easily fill out this chapter by describing each of the more
important games in greater detail.}
\VUsaut{6.0mm}
\VUfilet{22mm}
